%%%%%%%%%%%%%%%%%%%%%%%%%%%%%%%%%%%%%%%%%%%%%%%%%%%%%%%%%%%%%%%%%%%%%%%%
%%%%%%%%%%%%%%%%%%%%%% Simple LaTeX CV Template %%%%%%%%%%%%%%%%%%%%%%%%
%%%%%%%%%%%%%%%%%%%%%%%%%%%%%%%%%%%%%%%%%%%%%%%%%%%%%%%%%%%%%%%%%%%%%%%%

%%%%%%%%%%%%%%%%%%%%%%%%%%%%%%%%%%%%%%%%%%%%%%%%%%%%%%%%%%%%%%%%%%%%%%%%
%% NOTE: If you find that it says                                     %%
%%                                                                    %%
%%                           1 of ??                                  %%
%%                                                                    %%
%% at the bottom of your first page, this means that the AUX file     %%
%% was not available when you ran LaTeX on this source. Simply RERUN  %%
%% LaTeX to get the ``??'' replaced with the number of the last page  %%
%% of the document. The AUX file will be generated on the first run   %%
%% of LaTeX and used on the second run to fill in all of the          %%
%% references.                                                        %%
%%%%%%%%%%%%%%%%%%%%%%%%%%%%%%%%%%%%%%%%%%%%%%%%%%%%%%%%%%%%%%%%%%%%%%%%

%%%%%%%%%%%%%%%%%%%%%%%%%%%% Document Setup %%%%%%%%%%%%%%%%%%%%%%%%%%%%

% Don't like 10pt? Try 11pt or 12pt
\documentclass[10pt]{article}

\usepackage[utf8]{inputenc}
\usepackage[T1]{fontenc}

% LaTeX will typeset using Computer Modern Roman, which a lot of
% non-mathematicians and non-engineers won't like. Also, a few PDF
% viewers may not render CMR very well. Instead, Times New Roman can
% be used. That's what this package does.
\usepackage{times}
% \usepackage{libertinus}
% \usepackage{tgpagella}
% The automated optical recognition software used to digitize resume
% information works best with fonts that do not have serifs. This
% command uses a sans serif font throughout. Uncomment both lines (or at
% least the second) to restore a Roman font (i.e., a font with serifs).
% (NOTE: This requires the times package above)
%\renewcommand{\familydefault}{\sfdefault}

% This is a helpful package that puts math inside length specifications
\usepackage{calc}

% This package helps LaTeX auto-hyphenate hyphenated words if you use
% special hyphens. For example, bio\-/mimicry will properly hyphenate
% ``mimicry'' if necessary.
\usepackage[shortcuts]{extdash}

% Layout: Puts the section titles on left side of page
\reversemarginpar

%
%         PAPER SIZE, PAGE NUMBER, AND DOCUMENT LAYOUT NOTES:
%
% The next \usepackage line changes the layout for CV style section
% headings as marginal notes. It also sets up the paper size as either
% letter or A4. By default, letter was used. If A4 paper is desired,
% comment out the letterpaper lines and uncomment the a4paper lines.
%
% As you can see, the margin widths and section title widths can be
% easily adjusted.
%
% ALSO: Notice that the includefoot option can be commented OUT in order
% to put the PAGE NUMBER *IN* the bottom margin. This will make the
% effective text area larger.
%
% IF YOU WISH TO REMOVE THE ``of LASTPAGE'' next to each page number,
% see the note about the +LP and -LP lines below. Comment out the +LP
% and uncomment the -LP.
%
% IF YOU WISH TO REMOVE PAGE NUMBERS, be sure that the includefoot line
% is uncommented and ALSO uncomment the \pagestyle{empty} a few lines
% below.
%

\usepackage[paper=letterpaper,
            marginparwidth=1.30in,
            marginparsep=.01in,
            margin=0.75in,
            includemp]{geometry}

% Shift the main/details column 0.2 inches to the right
\addtolength{\textwidth}{-0.2in}
\addtolength{\oddsidemargin}{0.2in}
\addtolength{\evensidemargin}{0.4in}

%% Use these lines for A4-sized paper
%\usepackage[paper=a4paper,
%            %includefoot, % Uncomment to put page number above margin
%            marginparwidth=30.5mm,    % Length of section titles
%            marginparsep=1.5mm,       % Space between titles and text
%            margin=25mm,              % 25mm margins
%            includemp]{geometry}

%% More layout: Get rid of indenting throughout entire document
\setlength{\parindent}{0in}

% Provides special list environments and macros to create new ones
\usepackage[shortlabels]{enumitem}

% Simpler bibsections for CV sections
% (thanks to natbib for inspiration)
%
% * For lists of references with hanging indents and no numbers:
%
%   \begin{bibsection}
%       \item ...
%   \end{bibsection}
%
% * For numbered lists of references (with hanging indents):
%
%   \begin{bibenum}
%       \item ...
%   \end{bibenum}
%
%   Note that bibenum numbers continuously throughout. To reset the
%   counter, use
%
%   \restartlist{bibenum}
%
%   at the place where you want the numbering to reset.


\usepackage{graphicx}


\makeatletter
\newlength{\bibhang}
\setlength{\bibhang}{5em}
\newlength{\bibsep}
 {\@listi \global\bibsep\itemsep \global\advance\bibsep by\parsep}
\newlist{bibsection}{itemize}{3}
\setlist[bibsection]{label=,leftmargin=\bibhang,%
        itemindent=-\bibhang,
        itemsep=\bibsep,parsep=\z@,partopsep=0pt,
        topsep=0pt}
\newlist{bibenum}{enumerate}{3}
\setlist[bibenum]{label=[\arabic*],resume,leftmargin={\bibhang+\widthof{[999]}},%
        itemindent=-\bibhang,
        itemsep=\bibsep,parsep=\z@,partopsep=0pt,
        topsep=0pt}
\let\oldendbibenum\endbibenum
\def\endbibenum{\oldendbibenum\vspace{-.6\baselineskip}}
\let\oldendbibsection\endbibsection
\def\endbibsection{\oldendbibsection\vspace{-.6\baselineskip}}
\makeatother

%% Reference the last page in the page number
%
% NOTE: comment the +LP line and uncomment the -LP line to have page
%       numbers without the ``of ##'' last page reference)
%
% NOTE: uncomment the \pagestyle{empty} line to get rid of all page
%       numbers (make sure includefoot is commented out above)
%
\usepackage{fancyhdr,lastpage}
\pagestyle{fancy}
%\pagestyle{empty}      % Uncomment this to get rid of page numbers
\fancyhf{}\renewcommand{\headrulewidth}{0pt}
\fancyfootoffset{\marginparsep+\marginparwidth}
\newlength{\footpageshift}
\setlength{\footpageshift}
          {0.5\textwidth+0.5\marginparsep+0.5\marginparwidth-2in}
\lfoot{\hspace{\footpageshift}%
       \parbox{4in}{\, \hfill %
                    \arabic{page} of \protect\pageref*{LastPage} % +LP
%                    \arabic{page}                               % -LP
                    \hfill \,}}

% Finally, give us PDF bookmarks
\usepackage{color,hyperref}
\definecolor{darkblue}{rgb}{0.0,0.0,0.3}
\hypersetup{colorlinks,breaklinks,
            linkcolor=darkblue,urlcolor=darkblue,
            anchorcolor=darkblue,citecolor=darkblue}

%%%%%%%%%%%%%%%%%%%%%%%% End Document Setup %%%%%%%%%%%%%%%%%%%%%%%%%%%%


%%%%%%%%%%%%%%%%%%%%%%%%%%% Helper Commands %%%%%%%%%%%%%%%%%%%%%%%%%%%%

%%% HEADING AT TOP OF CURRICULUM VITAE

% The title (name) with a horizontal rule under it
% (optional argument typesets an object right-justified across from name
%  as well)
%
% Usage: \makeheading{name}
%        OR
%        \makeheading[right_object]{name}
%
% Place at top of document. It should be the first thing.
% If ``right_object'' is provided in the square-braced optional
% argument, it will be right justified on the same line as ``name'' at
% the top of the CV. For example:
%
%       \makeheading[\emph{Curriculum vitae}]{Your Name}
%
% will put an emphasized ``Curriculum vitae'' at the top of the document
% as a title. Likewise, a picture could be included:
%
%   \makeheading[{\includegraphics[height=1.5in]{my_picture}}]{Your Name}
%
% the picture will be flush right across from the name. For this example
% to work, make sure the extra set of curly braces is included. Also
% makes ure that \usepackage{graphicx} is somewhere in the preamble.
\newcommand{\makeheading}[2][]%
        {\hspace*{-\marginparsep minus \marginparwidth}%
         \begin{minipage}[t]{\textwidth+\marginparwidth+\marginparsep}%
             {\large \bfseries #2 \hfill #1}\\[-0.15\baselineskip]%
                 \rule{\columnwidth}{3pt}%
         \end{minipage}}

%%% SECTION HEADINGS

% The section headings. Flush left in small caps down pseudo-margin.
%
% Usage: \section{section name}
\renewcommand{\section}[1]{\pagebreak[3]%
    \vspace{1.5\baselineskip}%
    \phantomsection\addcontentsline{toc}{section}{#1}%
    \noindent\llap{\sffamily\bfseries\smash{\parbox[t]{\marginparwidth}{%
        \hyphenpenalty=10000\raggedright\MakeUppercase{#1}}}}%
    \vspace{-\baselineskip}\par}
%%% LISTS

% This macro alters a list by removing some of the space that follows the list
% (is used by lists below)
\newcommand*\fixendlist[1]{%
    \expandafter\let\csname preFixEndListend#1\expandafter\endcsname\csname end#1\endcsname
    \expandafter\def\csname end#1\endcsname{\csname preFixEndListend#1\endcsname\vspace{-0.6\baselineskip}}}

% These macros help ensure that items in outer-type lists do not get
% separated from the next line by a page break
% (they are used by lists below)
\let\originalItem\item
\newcommand*\fixouterlist[1]{%
    \expandafter\let\csname preFixOuterList#1\expandafter\endcsname\csname #1\endcsname
    \expandafter\def\csname #1\endcsname{\let\oldItem\item\def\item{\pagebreak[2]\oldItem}\csname preFixOuterList#1\endcsname}
    \expandafter\let\csname preFixOuterListend#1\expandafter\endcsname\csname end#1\endcsname
    \expandafter\def\csname end#1\endcsname{\let\item\oldItem\csname preFixOuterListend#1\endcsname}}
\newcommand*\fixinnerlist[1]{%
    \expandafter\let\csname preFixInnerList#1\expandafter\endcsname\csname #1\endcsname
    \expandafter\def\csname #1\endcsname{\let\oldItem\item\let\item\originalItem\csname preFixInnerList#1\endcsname}
    \expandafter\let\csname preFixInnerListend#1\expandafter\endcsname\csname end#1\endcsname
    \expandafter\def\csname end#1\endcsname{\csname preFixInnerListend#1\endcsname\let\item\oldItem}}

% An itemize-style list with lots of space between items
%
% Usage:
%   \begin{outerlist}
%       \item ...    % (or \item[] for no bullet)
%   \end{outerlist}
\newlist{outerlist}{itemize}{3}
    \setlist[outerlist]{label=\enskip\textbullet,leftmargin=*,topsep=0.5em,partopsep=0pt}
    \fixendlist{outerlist}
    \fixouterlist{outerlist}

% An environment IDENTICAL to outerlist that has better pre-list spacing
% when used as the first thing in a \section
%
% Usage:
%   \begin{lonelist}
%       \item ...    % (or \item[] for no bullet)
%   \end{lonelist}
\newlist{lonelist}{itemize}{3}
    \setlist[lonelist]{label=\enskip\textbullet,leftmargin=*,partopsep=0pt,topsep=0em}
    \fixendlist{lonelist}
    \fixouterlist{lonelist}

% An itemize-style list with little space between items
%
% Usage:
%   \begin{innerlist}
%       \item ...    % (or \item[] for no bullet)
%   \end{innerlist}
\newlist{innerlist}{itemize}{3}
    \setlist[innerlist]{label=\enskip\textbullet,leftmargin=*,parsep=0pt,itemsep=0pt,topsep=-0.5em,partopsep=-5pt}
    \fixinnerlist{innerlist}

% An environment IDENTICAL to innerlist that has better pre-list spacing
% when used as the first thing in a \section
%
% Usage:
%   \begin{loneinnerlist}
%       \item ...    % (or \item[] for no bullet)
%   \end{loneinnerlist}
\newlist{loneinnerlist}{itemize}{3}
    \setlist[loneinnerlist]{label=\enskip\textbullet,leftmargin=*,parsep=-1pt,itemsep=0pt,topsep=-1em,partopsep=0pt}
    \fixendlist{loneinnerlist}
    \fixinnerlist{loneinnerlist}

%%% EXTRA SPACE

% To add some paragraph space between lines.
% This also tells LaTeX to preferably break a page on one of these gaps
% if there is a needed pagebreak nearby.
\newcommand{\blankline}{\quad\pagebreak[3]}
\newcommand{\halfblankline}{\quad\vspace{-0\baselineskip}\pagebreak[3]}

%%% FORMATTING MACROS

% Provides a linked \doi{#1} that links doi:#1 to http://dx.doi.org/#1
\usepackage{doi}
% To change the text before the DOI, adjust this command
%\renewcommand\doitext{doi:}

% Provides a linked \url{#1} that doesn't require escape characters
\usepackage{url}

% You can adjust the style \url{} uses here:
% (options are: same, rm, sf, tt; defaults to tt)
\urlstyle{same}

% For \email{ADDRESS}, links ADDRESS to the url mailto:ADDRESS
% (uncomment to typeset the e\-/mail address in typewriter font;
%  otherwise, will be typeset in the \urlstyle above)
%\DeclareUrlCommand\emaillink{\urlstyle{tt}}
\providecommand*\emaillink[1]{\nolinkurl{#1}}
\providecommand*\email[1]{\href{mailto:#1}{\emaillink{#1}}}

\providecommand\BibTeX{{B\kern-.05em{\sc i\kern-.025em b}\kern-.08em \TeX}}
\providecommand\Matlab{\textsc{Matlab}}

% % Custom hyphenation rules for words that LaTeX has trouble with
% \hyphenation{bio-mim-ic-ry bio-in-spi-ra-tion re-us-a-ble pro-vid-er Media-Wiki}

%%%%%%%%%%%%%%%%%%%%%%%% End Helper Commands %%%%%%%%%%%%%%%%%%%%%%%%%%%

%%%%%%%%%%%%%%%%%%%%%%%%% Begin CV Document %%%%%%%%%%%%%%%%%%%%%%%%%%%%

\begin{document}
\makeheading{\LARGE Pramodit~Adhikari}

\section{Contact Information}

% NOTE: Mind where the & separators and \\ breaks are in the following
%       table. The table is one row made up of three parboxes. The left
%       parbox has address info, the middle parbox has a vertical bar,
%       and the right parbox has a phone and electronic contact
%       information.
%
% MACROS: \rcollength is the width of the right column of the table
%             (adjust it to your liking; default is 1.85in).
%         \spacewidth is the width of the area between the left and right boxes.
%
\newlength{\rcollength}\setlength{\rcollength}{2.20in}%
\newlength{\spacewidth}\setlength{\spacewidth}{20pt}
%
\begin{tabular}[t]{@{}p{\textwidth-\rcollength-\spacewidth}@{}p{\spacewidth}@{}p{\rcollength}}%

% Address box
\parbox{\textwidth-\rcollength-\spacewidth}{%
%Assistant Professor\\
\href{https://www.unl.edu/}{University of Nebraska-Lincoln}\\
\href{https://engineering.unl.edu/durhamschool/}{Durham School of Architectural Engineering and Construction}\\
Peter Kiewit Institute, \\
1110 S. 67th Street, Omaha, NE, USA}




&
% Uncomment to add a vertical bar in the middle of contact information
%{\vrule width 0.5pt}
\parbox[m][5\baselineskip]{\spacewidth}{} &

% Non-snail-mail contact information
\parbox{\rcollength}{%
\textit{Phone:} +1 720-491-9089 \\
\textit{E-~mail:} ~\href{mailto:adhikaripramodit@huskers.unl.edu}\email{adhikaripramodit@unl.edu}}
%\textit{website:} \href{https://therandomtalk.com/}{https://therandomtalk.com} (ongoing development)}
\end{tabular} 
%%
%% In modern CV's, it seems like ``Objective'' is frowned upon. Instead,
%% incorporate it into a well-constructed cover letter. The ``More
%% information'' can go at the end of the CV, but it should not distract
%% from the section giving references available to contact.
%%
%
% \section{Objective}
%
% Placement in an academic position (i.e., faculty, postdoctoral, or
% research scientist) that allows for advanced research in distributed
% complex adaptive systems (i.e., modeling, analysis, design, and
% verification) with a particular focus on the control of engineered
% agents (e.g., for communications, control, software, electronics, and
% sustainability) and the analysis of biological phenomena (e.g.,
% self-organization, ecological rationality)
% \begin{innerlist}
% \item More information and auxiliary documents can be found at\\\url{http://www.tedpavlic.com/facjobsearch/}
% \end{innerlist}


\section{Education}

\href{https://www.unl.edu/}{\textbf{University of Nebraska-Lincoln}}, Omaha, NE
\begin{outerlist}
	
	\item[] Ph.D., \href{https://engineering.unl.edu/durhamschool/}{Architectural Engineering}, Jan 2024 -- Present
	\begin{innerlist}
		\item Transferred from \href{https://engineering.unl.edu/cee/}{Civil and Environmental Engineering} (Jan 2021 -- Dec 2023)
		\item Adviser: \href{https://engineering.unl.edu/durhamschool/person/milad-roohi/}{Professor Milad Roohi}
		\item Area of Study: Structural Engineering
	\end{innerlist}
\end{outerlist}
\vspace{0.1in}

\href{https://www.colostate.edu/}{\textbf{Colorado State University}},
Fort Collins, CO
\begin{outerlist}
	\item[] M.Sc.,
	\href{https://engineering.unl.edu/cee/}
	{Department of Civil and Environmental Engineering}, Aug 2020
	\begin{innerlist}
		\item \textbf{Thesis} Topic: \emph{Life-Cycle Cost and Carbon-Footprint Analysis for Buildings and Communities Subjected to Tornadoes}
		\item Advisers:
		\begin{itemize}
			\item \href{https://www.engr.colostate.edu/ce/people/bruce-ellingwood/}
			{Professor Bruce R.~Ellingwood}
			\item \href{https://engineering.vanderbilt.edu/bio/hussam-mahmoud/} {Professor Hussam Mahmoud}
		\end{itemize}
		
		\item Area of Study: Structural Engineering
	\end{innerlist}
\end{outerlist}
\vspace{0.1in}

\href{https://tu.edu.np/}{\textbf{Tribhuvan University — Thapathali Campus}},
Thapathali, Kathmandu, Nepal
\begin{outerlist}
	\item[] B.S.,
	\href{https://tcioe.edu.np/}
	{Civil Engineering}, Jan 2017
	\begin{innerlist}
		\item Structural specialization (emphasis on earthquakes)
	\end{innerlist}
\end{outerlist}


\section{Current Academic Appointments}

\textbf{Graduate Research Assistant},
            \href{https://www.unl.edu/}{University of Nebraska-Lincoln}
            \hfill {Jan 2024 to present}
\begin{innerlist}

    \item[] \href{https://engineering.unl.edu/durhamschool/}{Durham School of Architectural Engineering and Construction} \\
    
    \begin{innerlist}
        \item \textit{Ongoing Research:}
            \begin{innerlist}
                \item Community Resilience against Tornado Hazards
                \item Bayesian Fragility Update and System Identification using Sparse Data
                \item Investigating Seismic Margin of Safety with Minimum Perturbation\\
            \end{innerlist}
    \end{innerlist}
\end{innerlist}


\section{Previous Academic Appointments}

\textbf{Graduate Research Assistant},
\href{https://www.unl.edu/}{University of Nebraska-Lincoln}
\hfill {Jan 2021 to Aug 2023}

\begin{innerlist}
	\item[] \href{https://engineering.unl.edu/cee/}{Department of Civil and Environmental Engineering}\\
	\begin{innerlist}
		\item \textit{Responsibilities:}
		\begin{innerlist}
			\item Cement-based electrochemical cells for large-scale energy storage
			\item Nondestructive diagnosis and probabilistic prognosis of aging plastic pipe\\
		\end{innerlist}       
	\end{innerlist}
\end{innerlist}



\textbf{Graduate Teaching Assistant},
            \href{https://www.unl.edu/}{University of Nebraska-Lincoln}
            \hfill {Jan 2021 to May 2021}

\begin{innerlist}
    \item[] \href{https://engineering.unl.edu/cee/}{Department of Civil and Environmental Engineering}\\
    \begin{innerlist}
        \item \textit{Responsibilities:}
            \begin{innerlist}
                \item Instructed students in the use of laboratory equipment for concrete mix design, tensile tests, compression tests, and torsion tests for the laboratory session of the required 3-credit course "CIVE 378 Materials of Construction"\\
            \end{innerlist}       
    \end{innerlist}
\end{innerlist}


\textbf{Graduate Research Assistant},
            \href{https://www.colostate.edu/}{Colorado State University}
            \hfill {Jan 2018 to Aug 2020}
\begin{innerlist}

    \item[] \href{https://www.engr.colostate.edu/ce/}{Department of Civil and Environmental Engineering} \\
    \begin{innerlist}
        \item \textit{Research:}
            \begin{innerlist}
                \item Developed tornado and hurricane fragility functions for 3 archetypes to be implemented in residential building and community-level economic and carbon-footprint analysis
                \item Modeled archetypes and residential communities, and created a database for cost and carbon footprint values of building components
                \item Established a generalized cost and carbon-footprint life-cycle estimation and analysis methodology for examining the benefits of different building practices for residential light-frame wood construction subjected to tornado hazards
                \item Probabilistic simulation of a tornado passing through the modeled community
                \item Developed wave/surge and hurricane fragility functions for archetypes to compare different retrofit strategies for a residential community near the shoreline and obtain an optimal strategy for damage mitigation \\
            \end{innerlist}
            \begin{innerlist}
                \item \textbf{Thesis: }\textit{Life-Cycle Cost and Carbon-Footprint Analysis for Buildings and Communities Subjected to Tornadoes}
            \end{innerlist}
    \end{innerlist}
\end{innerlist}


\section{Previous Professional Appointments}


\textbf{Wind-driven Fire Reliability Fellow}, \\
\href{ttps://ibhs.org/}{Insurance Institute of Business and Home Safety}
\hfill {May 2025 to July 2025}

\begin{innerlist}
	\item[] \href{https://ibhs.org/risk-research/wildfire/}{Wildfire Lab, and SC Wind \& Hail Lab} \\
	\begin{innerlist}
		\item \textit{Responsibilities:}
		\begin{innerlist}
			\item Experimental Design for Fire Resistance of Building Components
			\begin{itemize}
				\item Collaboratively helped design an experimental methodology in a cone calorimeter to quantify the resistance of wood siding.
			\end{itemize}
			\item Surface Temperature Prediction of Wood Siding
			\begin{itemize}
				\item Ideated a methodology for surface temperature prediction using a filtering algorithm for wood siding samples tested in a cone calorimeter \\
			\end{itemize}
		\end{innerlist}       
	\end{innerlist}
\end{innerlist}


\textbf{Internship}, \\
\href{https://inl.gov/}{Idaho National Laboratory}
\hfill {Oct 2022 to Dec 2023}

\begin{innerlist}
	\item[] \href{https://inl.gov/integrated-energy/hydrogen/}{Hydrogen and Electrochemistry Department } \\
	\begin{innerlist}
		\item \textit{Responsibilities:}
		\begin{innerlist}
			\item Proton-electrochemical cells manufacturing and testing platform
			\begin{itemize}
				\item Collaboratively helped design and implement an electrochemical cell testing platform to enhance research capabilities and testing throughput.
			\end{itemize}
			\item Sealing Research and Development
			\begin{itemize}
				\item Conducted analysis and redesign of sealant configurations, to improve the efficiency of electrochemical cells by optimizing material and durability.\\
			\end{itemize}
		\end{innerlist}       
	\end{innerlist}
\end{innerlist}



\section{Research Interests}

Bayesian statistics; community and infrastructure resilience; system reliability; natural hazards; vulnerability analysis; energy and sustainability;  game theory; resource allocation; machine learning; engineering education


% \section{Submitted Journal Publications}
%
% % Add a little space to nudge next ``Ref'd Journal Publications'' marginpar
% % down to make room for tall ``Submitted Journal Publications''
% % marginpar. If there are enough submitted journal publications, this
% % space will not be needed (and should be removed).
% \vspace{0.1in}

\section{Refereed Journal Publications}

\begin{bibenum}
    \item Adhikari, P., Roohi, M., \& Hernandez, E.  "Minimal perturbation of seismic records - Application to nonlinear time history analysis of building structures." \emph{Engineering Structures}, \textbf{Under review}.

    \item Adhikari, P., Roohi, M., Wood, R.L., \& Roueche, D. "Assessing Predictive Accuracy of Tornado Resilience Models Through Post-event Reconnaissance Data Analysis." \emph{ASCE Natural Hazard}, \textbf{Under review}.

    \item Adhikari, P., Mahmoud, H., Xie, A., Simonen, K., \& Ellingwood, B. 
    “Life-cycle cost and carbon footprint analysis for light-framed residential buildings subjected to tornado hazard.” 
    \emph{Journal of Building Engineering}, 32, 101657.\\
    \doi{https://doi.org/10.1016/j.jobe.2020.101657}

    \item Adhikari, P., Abdelhafez, M. A., Dong, Y., Guo, Y., Mahmoud, H. N., \& Ellingwood, B. R. 
    “Achieving residential coastal communities resilient to tropical cyclones and climate change.” 
    \emph{Frontiers in Built Environment}, 6, 576403.\\
    \doi{https://doi.org/10.3389/fbuil.2020.576403}

    \item Adhikari, P.~, Mahmoud, H. N.~, \& Ellingwood, B. R.
    Life-cycle cost and sustainability analysis of light-frame wood residential communities exposed to tornadoes.”
    \emph{Natural Hazards}, 109, 523-544.\\\doi{https://doi.org/10.1007/s11069-021-04847-x}

\end{bibenum}
% Add a little space to nudge next ``Conference Publications'' marginpar
% down to make room for tall ``Submitted Conference Publications''
% marginpar. If there are enough submitted journal publications, this
% space will not be needed (and should be removed).
\vspace{0.1in}

\section{Conference Proceedings (Peer-Reviewed)}
\begin{bibenum}
    \item Adhikari, P., Mahmoud, H., Xie, A., Simonen, K., \& Ellingwood, B. “Life-cycle cost and carbon footprint analysis for light-framed residential buildings subjected to tornado hazard.” \emph{Residential Building Design and Construction Conference}, State College, PA, USA, March 2020.
    
\end{bibenum}
\vspace{0.1in}

\section{Conference Abstracts/\\Presentations}
\begin{bibenum}
    \item Adhikari, P. \& Roohi, M. “Field-Calibrated Reliability and Resilience Framework for Tornado Design of Residential Buildings.” EMI Conference 2026, University of Boulder, Colorado, Boulder, CO, June 2-6, 2026.\vspace{-2em}
    \begin{flushright}
    \textbf{Third Place in Student Paper Competition.}
    \end{flushright}
    \item Adhikari, P. \& Roohi, M. “Bayesian Methods for Tornado Resilience Model Updating with Empirical Measurements.” EMI Conference 2025, University of California, Irvine, Anaheim Marriott, Anaheim, CA, May 27–30, 2025.
    \item Adhikari, P. \& Roohi, M. “Leveraging Machine Learning and AI for Enhanced Tornado Resilience in Community Infrastructure.” NHERI Computation Symposium 2025, University of California, Los Angeles, CA, Feb 2025.
    \item Adhikari, P., Roohi, M., Ghasemi S., Wood, R.L., \& Roueche, D.B. “Enhancing predictive accuracy in tornado community resilience models through hindcasting and post-event reconnaissance data analysis.” Engineering Mechanics Institute Conference 2024, Chicago, IL, May 2024.
    \item Adhikari, P., Abdelhafez, M. A., Dong, Y., Guo, Y., Mahmoud, H. N., \& Ellingwood, B. R. “Achieving residential coastal communities resilient to tropical cyclones and climate change.” Mini-Symposium on Safety Assessment of Aging Infrastructure: From Data to Decision, Engineering Mechanics Institute Conference, Columbia University, New York, NY, May 2020.
\end{bibenum}
\vspace{0.1in}

\section{Graduate Symposium}
\begin{bibenum}
    \item Adhikari, P., \& Roohi, M. "Model Updating with Empirical Measurements for Tornado Resilience Framework." 2026 COE Graduate Student Symposium, University of Nebraska-Lincoln, Lincoln, NE, Feb 2026
    \item Adhikari, P., Roueche, D.B., Wood, R.L., \& Roohi, M. "Investigating transferability and generalization of tornado resilience models through testbed juxtaposition." 2025 COE Graduate Student Symposium, University of Nebraska-Lincoln, Lincoln, NE, Feb 2025
    \item Adhikari, P., Roohi, M., Ghasemi S., Wood, R.L., \& Roueche, D.B. “Enhancing Community Resilience Models against Torando Hazards through Post-Tornado Reconnaissance Data.” 2024 COE Graduate Student Symposium, University of Nebraska-Lincoln, Lincoln, NE, May 2024
\end{bibenum}
\vspace{0.1in}

%\section{Books in Preparation}
%
%\begin{bibenum}
%    \item Pavlic, T.P., B.W.~Andrews, K.M.~Passino, and T.A.~Waite.
%        \emph{Foraging Theory for Engineering}. In preparation.
%\end{bibenum}

%\section{Papers in Preparation}
%
%\begin{bibenum}
%    \item Valentini, G., N.~Mizumoto, T.P.~Pavlic, S.C.~Pratt, and
%        S.I.~Walker. Control of information flow with acknowledgment
%        tokens in a non-human organism.
%
%    \item Xiaohui Guo, M.R.~Lin, L.P.~Saldyt, T.P.~Pavlic, Y.~Kang,
%        J.H.~Fewell. Adaptive alarm signal propagation by ant colonies:
%        using a machine learning approach to track individual and colony
%        response.
%
%    \item Pavlic, T.P., and S.C. Pratt. The Economic Framework: Using
%        constrained optimization to unify the ideal free distribution,
%        the marginal value theorem, and the geometric framework of
%        nutrition.
%
%    \item Pavlic, T.P. Risk-sensitive foraging and the Sharpe ratio.
%\end{bibenum}

%\textbf{Awaiting Decision}
%
%\begin{bibenum}
%    \item Senior staff, ``A new multi-objective optimization framework
%        for investigating mechanisms of social resource allocation'',
%        NIH, NIGMS, 2015. Revision in progress.
%
%\end{bibenum}
%
%\blankline



%\blankline
%
%\textbf{Not Awarded}
%
%\begin{bibenum}
%
%    \item Senior staff, ``Informational architecture of collective
%        decision making by \emph{Temnothorax} ants'', NSF, POLS, 2013.
%        Not awarded.
%
%    \item Senior staff,
%        ``Biological stoichiometry of horizontal gene transfer and the
%        social dynamics of microbial communities'', Army Research
%        Office, 2013. Not awarded.
%
%    \item Senior staff,
%        ``Biologically-inspired strategies for collective transport and
%        construction by multi-robot systems'', NSF, RI, 2013. Not
%        awarded.
%
%    \item Co\-/PI,
%        ``An Ant Model System for the Study of Nutrient Balance in
%        Social\-/Insect Pollinators'', USDA,
%        NIFA\-/AFRI Foundational proposal, 2013. Not awarded.
%
%    \item Senior staff,
%        ``Cooperative LED Arrays for Preference\-/Adaptive Lighting in
%        Smart Buildings'',
%        NSF,
%        EFRI\-/SEED preliminary proposal, 2009. Not awarded.
%
%\end{bibenum}

%\section{Academic Service}
%
%\textbf{Arizona State University, Tempe, AZ}
%\begin{lonelist}
%
%    \item College of Liberal Arts and Sciences Research Operations Committee,
%        \emph{Center for Social Dynamics and Complexity Representative},
%        2016--present.
%
%    \item Biosocial Complexity Initiative Directorate,
%        \emph{Liaison for Cross-University Activities},
%        2016--present.
%
%    \item Engineering Management Undergraduate Program Committee,
%        \emph{Member},
%        2015--present.
%
%    \item The Biomimicry Center,
%        \emph{Associate Director of Research},
%        2015--present.
%
%    \item Committee for the Development of Biomimicry and
%        Bio\-/inspired Research and Education Initiatives at ASU,
%        \emph{Chairman}.
%        2013.
%
%    \item Interdisciplinary Complexity Science Student Organization,
%        \emph{Founding faculty co\-/adviser}.
%        2013.
%
%\end{lonelist}

\section{Advising and Mentoring}
\textbf{Undergraduate Students}

\begin{outerlist}
    
    \item \textbf{Christian Anderson}, August 2025 - Present
    \begin{itemize}
        \item Developing a Climate-Informed Decision Support Framework for Urban Infrastructure Under Extreme Heat: A City of Lincoln Pilot. \\
        - \textit{Recently Funded. Fall, 2026 - Spring, 2027.}
    \end{itemize}

    \item \textbf{Tyler Reida}, Jan 2025 - Present
    \begin{itemize}
        \item Funded UCARE Project: Creating a Predictive Modelling Methodology for the Recovery of Communities after a Tornadic Event.
        \item Tyler presented at the Student Days in Lincoln
        \item Tyler presented at the 2026 NHERI symposium 
        \item Data collection for recovery analysis at various locations in the 2024 Nebraska and Iowa Tornadoes, paths (major areas included - Elkhorn and Blair, Nebraska, and Minden, Iowa).
    \end{itemize}
    
    \item \textbf{Geoffrey Sanderson}, Jan 2024 - May 2025
    \begin{itemize}
        \item Utilizing seismic response from sensors for response estimation of the building.
        \item Winner from the College of Engineering at the University of Nebraska–Lincoln’s Spring 2025 Research Days poster presentation
    \end{itemize}
    \item \textbf{Jordan Taylor}, Jul 2024 - Nov 2024
    \begin{itemize}
        \item Data collection efforts at Elkhorn, NE, about the recovery efforts from the Elkhorn tornado in April 2024.
    \end{itemize}
    
\end{outerlist}

\blankline

\section{Teaching Experience}
\href{https://www.unl.edu/}{University of Nebraska - Lincoln}

\begin{outerlist}
\item[] \textit{Graduate Assistant} \hfill Spring 2025
    \begin{innerlist}
        \item CIVE-859 (Reliability of Structures)
        \begin{itemize}
            \item Delivered a lecture on Conditional Probabilities and Bayesian Theory.
        \end{itemize}
    \end{innerlist}

\item[] \textit{Graduate Asisstant} \hfill Spring 2024
    \begin{innerlist}
        \item AREN-8940 (Infrastructure and Community Resilience) - 2024
        \begin{itemize}
            \item Delivered three lectures on Tornadoes, Hurricanes, and Fragility Analysis - formation of hazards, their impacts, and risk assessment methodologies.
        \end{itemize}
    \end{innerlist}

\item[] \textit{Graduate Asisstant} \hfill Spring 2026
    \begin{innerlist}
        \item AREN-8940 (Infrastructure and Community Resilience) - 2026
        \begin{itemize}
            \item Delivered one lectures on Tornadoes, and Fragility Analysis - formation of hazards, their impacts, and risk assessment methodologies.
        \end{itemize}
    \end{innerlist}
\end{outerlist}

\halfblankline


\section{Key Competencies}

Analytical
\vspace{0.05in}
\begin{innerlist}
    \item Probabilistic Analysis 
    \item Hazard and Exposure Modeling
    \item System Reliability Analysis
    \item Natural Hazards Vulnerability Analysis
    \item Community Resilience Analysis
\end{innerlist}
\vspace{0.1in}


Computer Programming:
\vspace{0.05in}
\begin{innerlist}[noitemsep,topsep=0pt]
    \item \includegraphics[height=1em]{./figures/python_logo.png} \hspace{0.7em} Python
    \item \includegraphics[height=1em]{./figures/r_logo.png} \hspace{0.4em} R
    \item \includegraphics[height=1em]{./figures/matlab_logo.png} \hspace{0.5em} \Matlab 
    \item \includegraphics[height=1em]{./figures/julia_logo.png}
    \hspace{0.5em} Julia \\
\end{innerlist}
\vspace{0.1in}

Computer Softwares
\vspace{0.05in}
\begin{innerlist}
    \item \includegraphics[height=1em]{./figures/autocad_logo.png}  \hspace{0.5em} AutoCAD
    \item \includegraphics[height=1em]{./figures/fusion_logo.png} \hspace{0.6em} Fusion 360
    \item \includegraphics[height=1em]{./figures/ai_logo.png} \hspace{2.6em} Adobe Ilustrator
    \item \includegraphics[height=1em]{./figures/inkscape_logo.png} \hspace{2.7em} Inkscape
    \item Finite Element Analysis
    \begin{itemize}
        \item \includegraphics[height=1em]{./figures/sap_logo.png} \hspace{0.5em} SAP 2000 
        \item \includegraphics[ scale=2, height=1em]{./figures/etabs_logo.png} \hspace{1.9em} ETABS
        \item \includegraphics[ scale=2, height=1em]{./figures/opensees_logo.png} \hspace{0.5em} OpenSeesPy
    \end{itemize}
\end{innerlist}
\vspace{0.1in}

Numerical Analysis:
\vspace{0.05in}
\begin{innerlist}
    \item \Matlab, R, Maple, Mathematica
\end{innerlist}
\vspace{0.1in}


Python, \Matlab, and R skill set:
\vspace{0.05in}
\begin{innerlist}
    \item Linear algebra, Fourier transforms, Monte Carlo analysis,
        statistics, Bayesian Statistics, visualization
\end{innerlist}
\vspace{0.1in}

Desktop Editing and Productivity Software:
\vspace{0.05in}
\begin{innerlist}
    \item VSCode, RStudio
    \item \LaTeX{},
    \item Microsoft Office, LibreOffice, OnlyOffice.com, Google Docs
\end{innerlist}
\vspace{0.1in}

Operating Systems:
\vspace{0.05in}
\begin{innerlist}
    \item Microsoft Windows, and Linux (Debian and Fedora family)
\end{innerlist}
\vspace{0.1in}

% \section{Expertise}
% Mathematics:
% \begin{innerlist}
%     \item Applied Mathematics, Real and Complex Analysis, Measure
%         Theory, Differential Geometry, Game Theory, Graph Theory,
%         Combinatorics
% \end{innerlist}

% \halfblankline

% Control Theory and Engineering:
% \begin{innerlist}
%     \item Linear and Nonlinear Systems Theory, Feedback, Variable
%         Structure Systems and Sliding Modes, Distributed and Intelligent
%         Control, Dynamic Optimization, Biomimicry, Bioinspiration,
%         Hybrid and CyberPhysical Systems
% \end{innerlist}

% \halfblankline

% Communications and Signal Processing:
% %
% \begin{innerlist}
%     \item Probability, Random Variables, Stochastic Processes,
%         Information Theory, Estimation, Networks
% \end{innerlist}

% \halfblankline

% Computer Science and Engineering:
% %
% \begin{innerlist}
%     \item Model Checking (automated, distributed, hybrid,
%         probabilistic), Hybrid Automata, Software Verification,
%         Component\-/Based Reusable Software
% \end{innerlist}

% \halfblankline

% Natural and Social Sciences (Biology, Neuroscience, Psychology, Anthropology):
% %
% \begin{innerlist}
%     \item Behavioral Ecology, Foraging Theory, Altruism, Impulsiveness,
%         Evolution
% \end{innerlist}

\section{Awards}

\href{https://www.unl.edu/}{University of Nebraska-Lincoln}
\begin{innerlist}
    \item \textbf{Graduate Research Assistantship}, University of Nebraska-Lincoln, June 2021 - present
    \item \textbf{Structural Engineers Association of Nebraska (SEAoN) Fellowship} Fund, Durham School of Architectural Engineering and Construction, University of Nebraska-Lincoln,  2026-2027 academic year
    \item \textbf{Third Place} in Student Paper Competition, EMI Conference 2026, University of Boulder, Colorado, Boulder, CO, June 2-6, 2026 
    \item \textbf{Graduate Student Travel Award} in the amount of \$500 on April 2024
    \item \textbf{Graduate Teaching Assistantship}, University of Nebraska-Lincoln, Jan 2021 - May 2021
    \item \textbf{Graduate Research Assistantship}, Colorado State University, Jan 2018 - August 2021
\end{innerlist}

        
\section{Professional Memberships}
\begin{innerlist}
	\item \href{https://www.asce.org/}{American Society of Civil Engineers} (Student Member)
	
    \item \href{https://my.eeri.org/}{Earthquake Engineering Research Institute} (Member)
    % \item \href{https://www.agu.org/}{American Geophysical Union} (Member) was a previous member
	% \item \href{https://engineering.unl.edu/durhamschool/earthquake-engineering-research-institute-eeri/}{Earthquake Engineering Research Institute (EERI)} (Student Member)
\end{innerlist}


\section{References}

\href{https://engineering.unl.edu/durhamschool/person/milad-roohi/}{\textbf{Professor Milad Roohi}}
(e\-/mail:~\href{mailto:milad.roohi@unl.edu}{milad.roohi@unl.edu})
\begin{innerlist}
    \item Assistant Professor, Durham School of Architectural Engineering and Construction, University of Nebraska-Lincoln
    \item[$\star$] \emph{Dr.~Roohi is my current Ph.D. supervisor.}
\end{innerlist}
\vspace{0.1in}

\href{https://www.engr.colostate.edu/ce/people/bruce-ellingwood/}{\textbf{Professor Bruce R.~Ellingwood}}
(e\-/mail:~\href{mailto:bruce.ellingwood@colostate.edu}{bruce.ellingwood@colostate.edu})
\begin{innerlist}
    \item Senior Research Scientist/Scholar, Colorado State University
    \item[$\star$] \emph{Dr.~Ellingwood was my Master's adviser.}
\end{innerlist}
\vspace{0.1in}

\href{https://engineering.unl.edu/civil/person/richard-wood/}{\textbf{Professor Richard Wood}}
(e\-/mail:~\href{mailto:rwood@unl.edu}{rwood@unl.edu})
\begin{innerlist}
	\item Associate Professor, Civil \& Environmental Engineering, University of Nebraska-Lincoln
	\item[$\star$] \emph{Dr.~Wood is co-author in an upcoming publication}
\end{innerlist}
\vspace{0.1in}

\href{https://eng.auburn.edu/directory/dbr0011}{\textbf{Professor David Roueche}}
(e\-/mail:~\href{mailto:dbr0011@auburn.edu}{dbr0011@auburn.edu})
\begin{innerlist}
	\item Associate Professor, Civil \& Environmental Engineering, Auburn University
	\item[$\star$] \emph{Dr.~Roueche is co-author in an upcoming publication}
\end{innerlist}
\vspace{0.1in}

\href{https://engineering.vanderbilt.edu/bio/?pid=hussam-mahmoud}{\textbf{Professor Hussam Mahmoud}}
(e\-/mail:~\href{mailto:hussam.mahmoud@vanderbilt.edu}{hussam.mahmoud@vanderbilt.edu})
\begin{innerlist}
    \item Professor, Civil and Environmental Engineering, Vanderbilt University
    \item[$\star$] \emph{Dr.~Mahmoud was my Master's co-adviser at Colorado State University.}
\end{innerlist}
\vspace{0.1in}



\end{document}
